\documentclass[11pt]{article}
\usepackage[utf8]{inputenc}
\usepackage[spanish]{babel}
\usepackage{amsmath,amssymb,amsthm}
\usepackage{graphicx}
\usepackage[margin=2.5cm]{geometry}
\usepackage{enumitem}
\title{Semana 5 -- Materia}
\author{Tutorías -- Introducción al Cálculo}
\date{31/08--04/09/2026}
\begin{document}
\maketitle

\section*{Introducción}

Esta semana cerramos el bloque de inecuaciones de la Sección 2.6 estudiando dos tipos que combinan
todo lo visto hasta ahora (axiomas de orden, valor absoluto, intervalos e inecuaciones racionales):
inecuaciones que involucran \textbf{valor absoluto} e inecuaciones que involucran \textbf{raíces}.
En ambos casos la idea central es la misma: transformar la inecuación original en una equivalente
libre de valor absoluto o de raíces, \emph{sin perder ni agregar soluciones}, prestando especial
atención a los signos de las expresiones involucradas y al dominio en el que la inecuación tiene
sentido.

\section{Inecuaciones con valor absoluto}

Recordemos la Definición 2.7:
\[
|x| = \begin{cases} x & \text{si } x \ge 0, \\ -x & \text{si } x < 0. \end{cases}
\]

De esta definición se deduce un par de equivalencias muy útiles para resolver inecuaciones. Para
$a > 0$:
\[
|x| < a \quad \leftrightarrow \quad -a < x < a,
\]
\[
|x| > a \quad \leftrightarrow \quad (x > a \ \lor\ x < -a).
\]

También es fundamental notar que si $A$ y $B$ son dos expresiones cualesquiera, entonces $|A|$ y
$|B|$ son siempre \textbf{no negativas}. Esto permite, cuando ambos miembros de una inecuación son
no negativos, elevar al cuadrado sin alterar el conjunto solución: en efecto, si $u \ge 0$ y
$v \ge 0$, entonces $u \le v \leftrightarrow u^2 \le v^2$ (esto se sigue de las propiedades del
orden de la Sección 2.3). En particular, $|A| \le |B| \leftrightarrow A^2 \le B^2$, pues
$|A|^2 = A^2$ y $|B|^2 = B^2$.

\subsection*{Ejemplo 1 (inspirado en el Problema 2.6 del libro)}

Resolver la inecuación
\[
\frac{|2x-1|}{|x+4|} \le 2.
\]

\textbf{Solución.} La expresión está definida para $x \neq -4$. Como el lado izquierdo es un
cociente de valores absolutos, es no negativo; el lado derecho también lo es. Luego podemos elevar
al cuadrado y obtener una inecuación equivalente:
\[
\frac{(2x-1)^2}{(x+4)^2} \le 4 \quad \leftrightarrow \quad (2x-1)^2 \le 4(x+4)^2,
\]
donde la segunda equivalencia usa que $(x+4)^2 > 0$ pues $x \neq -4$. Desarrollando:
\[
4x^2 - 4x + 1 \le 4(x^2 + 8x + 16) = 4x^2 + 32x + 64
\]
\[
\leftrightarrow \quad -4x + 1 \le 32x + 64 \quad \leftrightarrow \quad -36x \le 63 \quad
\leftrightarrow \quad x \ge -\frac{7}{4}.
\]
Como $-4 \notin \left[-\dfrac{7}{4}, \infty\right[$, la restricción de dominio no elimina ninguna
solución. El conjunto solución es
\[
\left[-\frac{7}{4}, \infty\right[.
\]

\subsection*{Ejemplo 2 (usando la equivalencia $|A|<B \leftrightarrow -B<A<B$)}

Resolver la inecuación
\[
|x+1| < 2x - 1.
\]

\textbf{Solución.} Para que exista alguna posibilidad de solución, el lado derecho debe ser
positivo (pues el lado izquierdo es siempre $\ge 0$), es decir, necesitamos $2x - 1 > 0$. Bajo esa
condición, la definición de valor absoluto nos da la equivalencia
\[
|x+1| < 2x-1 \quad \leftrightarrow \quad -(2x-1) < x+1 < 2x-1.
\]
Separamos en dos inecuaciones:
\[
-(2x-1) < x+1 \ \leftrightarrow\ -2x+1 < x+1 \ \leftrightarrow\ -3x < 0 \ \leftrightarrow\ x > 0,
\]
\[
x+1 < 2x - 1 \ \leftrightarrow\ 2 < x \ \leftrightarrow\ x > 2.
\]
Intersectando con la condición $2x-1>0$ (es decir $x > 1/2$), y con $x>0$, $x>2$, obtenemos que la
condición más restrictiva es $x>2$. El conjunto solución es
\[
]2,\infty[.
\]

\section{Inecuaciones con raíces}

Al resolver inecuaciones que involucran raíces cuadradas hay que tener presentes dos cuidados
adicionales respecto de las inecuaciones racionales ya vistas:

\begin{enumerate}[label=(\alph*)]
    \item \textbf{Dominio:} una expresión $\sqrt{f(x)}$ sólo está definida cuando $f(x) \ge 0$.
    Cualquier solución debe satisfacer esta condición.
    \item \textbf{Signo al elevar al cuadrado:} como $\sqrt{f(x)} \ge 0$ siempre, al comparar
    $\sqrt{f(x)}$ con una expresión $g(x)$ hay que fijarse en el signo de $g(x)$ antes de elevar al
    cuadrado, exactamente como hicimos con el valor absoluto en la sección anterior.
\end{enumerate}

Concretamente, si queremos resolver $\sqrt{f(x)} \le g(x)$, necesitamos $f(x) \ge 0$ y $g(x) \ge 0$
(si $g(x) < 0$ no hay solución posible, pues el lado izquierdo es no negativo), y bajo esas
condiciones la inecuación es equivalente a $f(x) \le g(x)^2$.

\subsection*{Ejemplo 1: una inecuación con raíz explícita}

Resolver la inecuación
\[
\sqrt{2x+3} \le x.
\]

\textbf{Solución.} Para que la raíz esté definida necesitamos $2x+3 \ge 0$, es decir
$x \ge -\dfrac{3}{2}$. Además, como el lado izquierdo es $\ge 0$, cualquier solución debe cumplir
también $x \ge 0$. Bajo esta última condición (que ya implica la anterior), podemos elevar al
cuadrado:
\[
2x + 3 \le x^2 \quad \leftrightarrow \quad x^2 - 2x - 3 \ge 0 \quad \leftrightarrow \quad
(x-3)(x+1) \ge 0 \quad \leftrightarrow \quad (x \ge 3 \ \lor\ x \le -1).
\]
Intersectando con la condición $x \ge 0$, sólo sobrevive $x \ge 3$. El conjunto solución es
\[
[3, \infty[.
\]

\subsection*{Ejemplo 2 (inspirado en el Problema 2.4 del libro): tabla de signos con raíces}

Resolver la inecuación
\[
\frac{(x^2-2)(x^2-7)}{(x-1)(x-6)} > 0.
\]

\textbf{Solución.} Notemos primero que esta es una inecuación racional cuyos factores cuadráticos
tienen raíces irracionales: $x^2 - 2 = 0 \leftrightarrow x = \pm\sqrt{2}$ y
$x^2 - 7 = 0 \leftrightarrow x = \pm\sqrt{7}$. Factorizando completamente,
\[
\frac{(x-\sqrt{2})(x+\sqrt{2})(x-\sqrt{7})(x+\sqrt{7})}{(x-1)(x-6)} > 0,
\]
con la restricción $x \neq 1$ y $x \neq 6$. El signo de esta expresión depende del signo de cada
factor, así que ordenamos los seis puntos críticos de menor a mayor. Como
$\sqrt{2} \approx 1.41$ y $\sqrt{7}\approx 2.65$, el orden es:
\[
-\sqrt{7} < -\sqrt{2} < 1 < \sqrt{2} < \sqrt{7} < 6.
\]
Llamemos $E(x)$ a la expresión completa. Estudiamos el signo de cada factor en cada intervalo:

\begin{center}
\renewcommand{\arraystretch}{1.3}
\begin{tabular}{|c|c|c|c|c|c|c|c|}
\hline
 & $\;]-\infty,-\sqrt{7}[\;$ & $-\sqrt{7}$ & $]-\sqrt{7},-\sqrt{2}[$ & $-\sqrt{2}$ & $]-\sqrt{2},1[$ & $1$ & $]1,\sqrt{2}[$ \\
\hline
$x-\sqrt{7}$   & $-$ & $-$ & $-$ & $-$ & $-$ & $-$ & $-$ \\
$x+\sqrt{7}$   & $-$ & $0$ & $+$ & $+$ & $+$ & $+$ & $+$ \\
$x-\sqrt{2}$   & $-$ & $-$ & $-$ & $-$ & $-$ & $-$ & $-$ \\
$x+\sqrt{2}$   & $-$ & $-$ & $-$ & $0$ & $+$ & $+$ & $+$ \\
$x-1$          & $-$ & $-$ & $-$ & $-$ & $-$ & $0$ & $+$ \\
$x-6$          & $-$ & $-$ & $-$ & $-$ & $-$ & $-$ & $-$ \\
\hline
$E(x)$         & $+$ & $0$ & $-$ & $0$ & $+$ & $*$ & $-$ \\
\hline
\end{tabular}
\end{center}

\begin{center}
\renewcommand{\arraystretch}{1.3}
\begin{tabular}{|c|c|c|c|c|c|c|}
\hline
 & $\sqrt{2}$ & $]\sqrt{2},\sqrt{7}[$ & $\sqrt{7}$ & $]\sqrt{7},6[$ & $6$ & $]6,\infty[$ \\
\hline
$x-\sqrt{7}$   & $-$ & $-$ & $0$ & $+$ & $+$ & $+$ \\
$x+\sqrt{7}$   & $+$ & $+$ & $+$ & $+$ & $+$ & $+$ \\
$x-\sqrt{2}$   & $0$ & $+$ & $+$ & $+$ & $+$ & $+$ \\
$x+\sqrt{2}$   & $+$ & $+$ & $+$ & $+$ & $+$ & $+$ \\
$x-1$          & $+$ & $+$ & $+$ & $+$ & $+$ & $+$ \\
$x-6$          & $-$ & $-$ & $-$ & $-$ & $0$ & $+$ \\
\hline
$E(x)$         & $0$ & $+$ & $0$ & $-$ & $*$ & $+$ \\
\hline
\end{tabular}
\end{center}

(el símbolo $*$ indica que $E(x)$ no está definida). Buscamos donde $E(x) > 0$, por lo que el
conjunto solución es
\[
]-\infty,-\sqrt{7}[\ \cup\ ]-\sqrt{2},1[\ \cup\ ]\sqrt{2},\sqrt{7}[\ \cup\ ]6,\infty[.
\]

\end{document}
